% !TeX program = xelatex
% ===========================================================================
%  对于音乐“非量化”工作的错位与反思
%   —— 与中国期刊投稿体例对齐；版式由 musepaper.sty 统一提供
%
%  编译：xelatex 三遍（compile.bat 已封装）
%  字体：共享目录 ../fonts/（思源宋体静态实例 400/700）
% ===========================================================================
\documentclass[11pt,a4paper,fontset=none]{ctexart}

% 共享字体目录（必须在 musepaper 之前声明）
\providecommand{\museFontPath}{../fonts/}
\usepackage{musepaper}

% 页眉（中文期刊体例：居中排题名）
\页眉{对于音乐“非量化”工作的错位与反思}

\title{对于音乐“非量化”工作的错位与反思}
\author{王子健}
\date{}

%% ---------------------------------------------------------------------------
%%  引用体例说明（《中央音乐学院学报》著录规范）：
%%    本刊所引文献统一以页下脚注体现，文后不再附录参考文献。
%%    故原 \begin{thebibliography} 已改为各处 \footnote。
%%    ⚠️ 规范要求著录「所引页码」；凡标【待核】/【X】处请按所据版本补齐。
%% ---------------------------------------------------------------------------
\begin{document}

% 题名页不加页眉与页码（期刊投稿惯例）；页码从第 2 页起自然流动。
\thispagestyle{empty}

% ===========================================================================
% 题名页
% ===========================================================================
\begin{center}
  {\sffamily\bfseries\fontsize{17}{21}\selectfont
   对于音乐“非量化”工作的错位与反思\par}
  \vspace{0.9em}
  {\normalsize 王子健\par}
  \vspace{0.35em}
  {\small（MozunStudio，安徽 合肥 230000）\par}
  \vspace{0.25em}
  {\small 通信邮箱：368164331@qq.com\par}
\end{center}
\vspace{0.4em}

\begin{摘要}
“音乐不可量化”长期被视为音乐行业的默认前提，音乐亦被习惯性地归入“人性的艺术”“心灵的感悟”。
本文认为，这一命题首先应当被精确化：究竟是“音乐”不可量化，还是“音乐情感”不可量化？
在界定量化为“离散化与符号化的有损压缩”的基础上，本文提出可量化之物的四项判据，
其中“对象×用途×容差”的三元关系构成核心；据此论证音乐存在“静止与运动”两个层面——
形式与结构层面静止且可量化，审美与人为相对动态层面不可量化，二者以“伴奏带与人声”式的协作模式共存。
本文进一步重审音乐定义，提出音乐之“规律”即“在一定时间内所占据的声音比例”，
确立时间维度在音乐中的本体地位；援引嵇康《声无哀乐论》与汉斯立克自律论，佐证审美层的人为属性；
并以朱载堉十二平均律、巴赫《平均律钢琴曲集》、记谱法演化史及作曲机制（动机与训练的同构性）
论证作曲编曲的可量化性。最后，本文揭示量化意识缺位的制度性根源——口传心授的传承模式、
“对错”评价的错位与学科建制的失衡，并展望“结构层绝对坐标、审美层显式容差”的
分层量化方案之于 AI 音乐编曲的意义。
\end{摘要}

\关键词{音乐量化；声音本体论；时间结构；音乐美学；分层量化；人工智能作曲}

\中图分类号{J601}
\文献标识码{A}

% ---------------- 英文块 ----------------
\英文题名{Dislocation and Reflection on the ``Non-Quantifiable'' Work of Music:\\
Quantifiability, Temporal Ontology and Layered Quantization}
\英文作者{WANG Zijian}
\英文单位{(MozunStudio, Hefei 230000, Anhui, China)}

\begin{英文摘要}
The proposition that ``music is non-quantifiable'' has long been regarded as a self-evident
premise in the music industry, and music has been habitually classified as ``the art of humanity''
and ``the perception of the soul''. This paper argues that this proposition should first be
refined: is it ``music'' that cannot be quantified, or ``musical emotion'' that cannot be
quantified? On the basis of defining quantization as ``lossy compression through discretization
and symbolization'', this paper proposes four criteria for quantifiable objects, of which the
ternary relation ``object $\times$ purpose $\times$ tolerance'' constitutes the core. Accordingly,
it demonstrates that music comprises two layers---the static and the dynamic: the layer of form
and structure is static and quantifiable, while the layer of aesthetic appreciation and
human-relative dynamics is non-quantifiable; the two coexist in a collaborative mode exemplified
by ``accompaniment track and human voice''. This paper further re-examines the definition of
music, proposing that the ``regularity'' of music is ``the proportion of sound occupying a given
period of time'', thereby establishing the ontological status of the temporal dimension in music.
Citing Ji Kang's \textit{Sheng Wu Ai Yue Lun} (On the Absence of Sadness and Joy in Sound) and
Hanslick's formalism, it corroborates the human attribute of the aesthetic layer; and, with Zhu
Zaiyu's twelve-tone equal temperament, Bach's \textit{The Well-Tempered Clavier}, the evolution of
musical notation, and the composition mechanism (the isomorphism of motif and training), it argues
for the quantifiability of composition and arrangement. Finally, this paper reveals the
institutional roots of the absence of quantification consciousness---the oral transmission model,
the dislocation of ``right/wrong'' evaluation, and the imbalance of disciplinary construction---and
looks forward to the significance of a layered quantization scheme (absolute coordinates at the
structural layer, explicit tolerance at the aesthetic layer) for AI music arrangement.
\end{英文摘要}

\英文关键词{music quantification; ontology of sound; temporal structure;
aesthetics of music; layered quantization; AI music composition}

\newpage
% ===========================================================================
% 目录
% ===========================================================================
{\sffamily\bfseries\fontsize{14}{17}\selectfont
 \begin{center}目\quad 录\end{center}}
\vspace{0.6em}
\renewcommand{\contentsname}{}
\vspace{-3.2em}
\tableofcontents
\newpage

% ===========================================================================
% 正文
% ===========================================================================
\section{引言：问题的提出}

在传统音乐观念中，“音乐不可量化”几乎是不言自明的公理。我们习惯把音乐归入“人性的艺术”
“心灵的感悟”，仿佛一旦动用尺度与网格，音乐便不再是音乐。这一定见如此牢固，以至于整个
行业的数字化进程——从打谱软件到数字音频工作站（DAW）——都不得不在“艺术性”的辩护下
小心翼翼地推进。

然而，笔者不揣浅陋，以为这一公理本身值得追问：我们所说的“音乐不可量化”，究竟是“音乐”
不可量化，还是“音乐情感”不可量化？这两个命题指向完全不同的对象：前者关乎声音组织的本体
属性，后者关乎聆听主体的心理效应。将二者混为一谈，正是音乐“非量化”工作中最根本的错位。

本文的论证路径如下：先界定“量化”及可量化之物的判据；再论证音乐存在可量化与不可量化的
分层，且二者协作共存；进而重审音乐的定义，确立时间维度的本体地位；随后以音乐史与创作
实践为证，考察作曲编曲的可量化性；最后揭示量化意识缺位的制度性根源，并展望分层量化的
可行方案。

\section{何为量化：概念界定与判据}

\subsection{量化的定义与结构}

从专业角度而言，量化＝把连续（或高基数）的对象集合，映射到一个离散（可数）的符号集合的
过程。通常，它由三个结构构件构成：其一，源集合——连续域或高分辨率域，如一段连续的时间、
一根连续的音高轴；其二，划分（partition）——把源集合切成互不相交的单元，每个单元称为
一个“桶”（bin）；其三，映射规则——每个桶指派一个代表符号（通常为其中心或质心），此后
所有落入该桶的元素，一律以该符号指称。

因此，量化从来不是一件事，而是两件事的合谋：离散化（切）加上符号化（命名）。切完了还要
能叫得出名字，量化才算完成。而它的本质，是一种“有损压缩”——把不可数压成可数，代价是
桶内差异被抹平。事实上，人类认知本身就是一场大量化：感官是连续的流，概念是切的桶，
“红”这个字，不过是千万种红色光谱的质心符号。

\subsection{可量化之物的判据}

可被量化的东西，必须同时满足四项判据：其一，存在差异结构——对象之间存在可比的距离或
等同关系，连“像不像”都无从谈起的东西，切无可切；其二，存在拓扑/相邻性——划分必须尊重
“邻近”，桶内之物应比桶间之物更为相似，否则划分便是乱砍；其三，存在可容忍的容差——桶内
差异对既定用途无关紧要。这一条最为关键：可量化性从来不是对象自身的属性，而是“对象×
用途×容差”三者的关系；其四，划分稳定可复现——同一对象量化十次，须得到同一个符号，
映射规则须是确定的函数。

据此，严格地说，并不存在“天然可量化的东西”，只存在“在某个容差下、为某个用途、可以被
有限划分的对象”。

\subsection{不可量化之物的病灶}

反观不可量化之物，其病灶同样清晰，约有三端：其一，差异在所有尺度上都同样致命——任何
容差都不能安全使用。芝诺的飞矢悖论即其典型：把运动切成瞬间，运动当场消失，运动的“活”
恰恰在连续性本身；其二，意义活在整体里，不在局部——切开即死，如乐句的呼吸、一句诗的
韵味；其三，没有稳定同一性——每次出现都不一样，无法比较，亦无法归类。

\subsection{不可量化论的四种归位}

§2.3 所列三端，回答的是“何物不可量化”；此处要处理的是另一个问题：主张“音乐
不可量化”的几种常见论证，各自“成立到哪一层为止”。二者相表里，结论却不同——
前者说的是对象的性质，后者说的是论证的边界。把它们逐一归位，本文的分层主张
也就有了入口。

\textbf{其一，连续统论证。}此说以为：量化的前提是划分，而音乐中最要紧的东西——
滑音、渐强、弹性速度——都是连续统；一旦切成离散的单元，连续性本身便告消失。
此论成立的部分是确凿的：连续统不能被“无损”离散化。但它随即越了界——“会丢
东西”并不等于“不能切”，只说明必须追问：丢掉的究竟是什么，它是否落在既定用途
的容差之内。这一追问正是判据三的实质。事实上，人类认知本身就是一场大量化：
感官所接的是连续光谱，而“红”这个字，不过是千万种红色被压成的一个符号。倘若
连续统论证可以成立，则连“红”这个字也不该存在。

\textbf{其二，整体论论证。}此说以为：音乐的意义活在整个乐句乃至整部作品之中，
切开即死，如乐句的呼吸、一段旋律的韵味。此论把“意义的归属”与“对象的存在”
混为一谈。一个乐句完全可以被完整地记为一组具有明确起止时刻的坐标——它的每一个
音都有确定的位置，而这并不妨碍它作为一个整体被听者把握；韵味是听者面对这组坐标
时产生的效应，它属于接收的一端，不属于声音组织本身。整体性并不排除可坐标化，
它只对坐标系统提出一个要求：必须能够容纳“关系”，而不只是彼此孤立的点。

\textbf{其三，主观性论证。}此说以为：音乐所引发的情感是私人的、不可通约的，
因而无从量化。此论恰恰是本文开篇所指那个错位的标本——它谈的是“听者的心理”，
却拿来给“音乐”下判断。情感是否可量化是另一个问题（心理学以效价与唤醒度两个
维度施测，本身即是一种量化尝试），但无论其答案如何，都不能反推声音组织不可测。
退一步说，即便承认情感世界拒绝刻度，那也只是为审美层划界，与结构层无涉。

\textbf{其四，历史性论证。}此说以为：音乐的价值随时代流转，昨日的不协和是今日
的常规，因而不存在稳定的量。此论反驳的是“审美标准的恒常性”，而非“结构的可测
性”。同一个大三和弦在 1722 年与在今日，其音高关系、时值配比、频谱结构并无二致；
变的是对它的“评价”，不是它本身。历史性非但不构成障碍，反而恰好印证了本文的
分层主张：结构层稳定，审美层流动。

综观四说，它们共享同一个结构：各自在“某一层”成立，却都被用来支撑一个“跨层”
的结论。这一越界之所以长期未被察觉，是因为“音乐”一词在讨论中始终没有被拆开——
它同时指涉声音组织与听觉经验，于是针对后者的论断便顺理成章地被读成了针对前者的
论断。把这四种论证各自归位到审美层，它们与结构层的可量化性便不再冲突；本文
第三章的分层，正是为了给它们各自找到该待的地方。

\section{音乐可量化性的辩证：静止与运动}

\subsection{对“演奏不可重复”论的反驳}

按照上述逻辑，或有人直接发问：音乐妥妥地是不可量化之物，你在讨论什么？的确，即便最优秀
的演奏家，也无法完成一模一样的两次演奏——这在物理上是不可能的。但这就证明音乐不可量化
吗？笔者以为未必。若音乐果真不可量化，又如何解释 DAW 与任意打谱软件对音符的识别与定位
功效？这是显然的事实，不能一概而论。用哲学的话来说，我们既要看到事物静止的一面，又要
看到事物运动的一面。

\subsection{音乐的分层：结构层与审美层}

那么，对音乐而言，什么是静止？什么是运动？笔者的结论是：就“音乐的形式与结构”而言，
音乐是静止且可量化的；就“音乐的审美”与“人为的相对动态”而言，音乐是不可量化的。

这并不难理解。假设你在 DAW 或打谱软件中记录下一个音符，下次播放，这个音符的位置一定
老老实实地待在那里；而若由人演奏，则可能出现微小的偏移——这种“偏移”恰恰是“人性化”的
体现。琴键与网格之间的游离，正是演奏者呼吸的痕迹。

\subsection{静止与运动的协作：以“伴奏带与人声”为例}

那么，二者之间一定水火不容吗？答案显然是否定的。凡唱过歌的人都有合伴奏带的经验：既然
伴奏带是静止的，人便会自觉地相对运动，以求保持相对平衡；若换作真人钢伴与演唱者之间，
也存在一种默契的平衡，否则整场表演便会直接乱套。由此可见，静止不是运动的敌人，而是
运动得以发生并保持有序的坐标系：正因为网格恒定，偏离才有意义；正因为坐标存在，呼吸
才有位置。静止与运动，各司其职，互为成全。

\section{音乐的本体重审：时间维度}

\subsection{通用音乐定义的审视}

乐理书中关于音乐的描述并无错误，但长期以来，我们忽略了一个关键的问题：什么是音乐？
通用的解释如“音乐是一种通过有规律的声波组织来表达思想、情感和社会生活的艺术形式与文化
活动”，又如“音乐是利用声音的高低、强弱、时长和音色等要素，通过节奏、旋律、和声等形式
构成的听觉艺术”。这些论述都对，但请注意其中的一句——“有规律的声波组织”：那么，这个
“规律”究竟是什么？而所谓“利用声音的高低、强弱、时长和音色等要素”，这一系列的“艺术”
又是什么？

\subsection{规律：一定时间内所占据的声音比例}

笔者的看法是：这个规律，即“在一定时间内所占据的声音比例”。这句话看似抽象，实则很好
理解：一场音乐会的时长总是相对固定的——两小时的内容，即便有微小误差，也不可能被压缩到
半小时，除非发生异常事故。而其中的音高也罢、休止也罢、不同乐器也罢，都不过是“时间内
可量化的子集”。

由此可以得出一个长期被忽视的结论：时间，才是音乐中真正具有本体地位的单位。一切声音
要素——音高（单位时间内的频率分布）、音色（随时间演化的频谱）、和声（同一时刻的纵向
切片）、节奏（声音在时间轴上的疏密分布）——皆可投影于时间轴之上。音乐不是声音的集合，
而是声音在时间中的组织。

\subsection{音乐本体与“人”的载体之辨}

或有人说，这种定义会打破音乐的“美学”与“艺术性”。笔者以为，这种担忧并不成立，二者
并不冲突，这是一个简单且长期被忽视的问题。我们不妨转换视角：现有的音源体系与 DAW 编曲
设备已然十分完善，那么编曲师以 DAW 做出的音乐作品，难道没有“美学”与“艺术性”吗？
显然不是。既然 DAW 本身就是一种音乐量化且为行业公认的工具，何来“打破美学艺术性”一说？
我们不能说用耳机听音乐不算欣赏，看音乐会才算欣赏——那显然不符合逻辑，因为“音乐的本体”
仍然是“声音”这一介质；至于它是否由人产生，只是“过程本身”，而非“最终结果”。

说得最直白的问题便是：观众花钱，究竟是“欣赏音乐”，还是“欣赏音乐的产生者”？如果只以
产生者来定义音乐，那么确实可以解释这一误区——所谓缺乏“美学”与“艺术性”，不过是因为
美学被绑定上了“人”这一载体而已。

但凡学习过哲学，便会熟知这样一句话：规律是不以人的载体为转移的；它可以被能动地利用，
但绝非因个人意志而转移。同理，既然音乐早被定义为“有规律的声波组织”，那么将“人”这一
载体强加于规律之上，是不是显得过于强势、违背了哲学本身？

\section{声无哀乐：审美层的人为属性}

\subsection{嵇康《声无哀乐论》的先声}

若从“量化的形式”出发，“音乐是心灵的体现”——这话乍一听没有问题，但它忽略了一点：
人类的心理，与音乐本身是无关的。这一思想，早在嵇康的美学著作《声无哀乐论》\footnote{〔三国魏〕嵇康：《声无哀乐论》，【待核：所据版本，如戴明扬《嵇康集校注》，北京：人民文学出版社，1962年】，第【X】页。}中
便已被清晰描述。嵇康认为“心之与声，明为二物”：声音是客观存在的自然音响，哀乐是精神
被触动后产生的感情，二者并无因果关系。也就是说，音乐本身不携带哀乐，哀乐存在于聆听者
那里。

\subsection{极值情境的思想实验}

我们默认音乐被欣赏的前提，是人处于一种无外力干扰且平和的状态中。但人类的心理活动本质
上是多元的，我们不能忽略极值这一因素。不妨做一个思想实验：假设一人中了 100 万元彩票，
另一人亏了 100 万元；那么，给亏钱者听“欢快的音乐”，是否真能让他感受到快乐？给赚钱者
听“悲伤的音乐”，是否当真使他悲伤？答案显而易见。

需要补充的是，此处所谓“无关”，更精确的表达是“不必然对应”：音乐可以诱发情绪反应，
但不能规定情绪反应——音乐是情绪的诱因，而非情绪的载体。这一措辞的严谨化，并不削弱
嵇康的立场，反而使其更有凭据。

\subsection{历史呼应：自律论与他律论之争}

在西方音乐美学史上，这一思想并非孤例。汉斯立克于 1854 年出版的《论音乐的美》\footnote{〔奥〕爱德华·汉斯立克：《论音乐的美——音乐美学的修改刍议》，杨业治译，北京：人民音乐出版社，【待核：出版年份】，第【X】页。}
提出著名的自律论：音乐的美不在于表现情感，而在于“乐音的运动形式”本身；与之相对，
浪漫主义的情感美学（他律论）则主张音乐的美在于表达情感。两派争执近两百年而未决。以本文
的“静止/运动”分层观之，二者的领地其实各得其所：结构层是自律的——形式即美，静止可
量化；审美层是他律的——因人而异，运动不可量化。本文无意调停二者的立场，而只是为二者
划定地界：从此不必再争，各归其位而已。

\section{作曲与编曲的可量化性考察}

\subsection{历史证据：从朱载堉到记谱法}

回到一个现实的问题：传统的作曲编曲已有如此悠久的历史，就“作曲与编曲”本身而言，是否是
可以量化的呢？答案仍然可以沿用前文的回答。

从历史的角度看：中国的朱载堉于万历十二年（1584 年）通过精密计算将八度音程均分为十二个
半音，每个半音的音程比值为 2 的 12 次方根（约 1.059463），实现了真正数学意义上的平均
律\footnote{〔明〕朱载堉：《律学新说》，【待核：点注者、出版社、出版年份】，第【X】页；《律吕精义》，【待核：同上或另行注明】，第【X】页。}——这一计算，是用八十一档大算盘开方完成的，精确到小数点后 24 位。
巴赫的《平均律钢琴曲集》（第一卷，1722 年）\footnote{〔德〕J. S. 巴赫：《平均律钢琴曲集》第一卷（BWV 846—869），【待核：所据版本，如 Leipzig: Bärenreiter, 1950】，第【X】页。}则从创作实践上印证了平均律的可行性。
至于记谱法，无论是中国传统的减字谱，还是西方的纽姆谱、四线谱乃至五线谱，都在为音乐的可
量化性做出努力。但笔者以为，我们不能将记谱法简单地归类为“为了保存、流传”，而更应回归
到量化与科学的本身——记谱，本质上就是音乐的时间结构与音高结构的符号化，即量化的第一次
实践。

\subsection{创作机制：动机、训练与生成}

从创作的角度看：作曲中确实不能排除一个因素——“动机”。在音乐中，我们常把“动机”归为
作曲家的“灵光乍现”，但灵光乍现真的是玄学吗？笔者以为不是。这与现代人工智能的
Transformer 机制\footnote{A. Vaswani, N. Shazeer, N. Parmar, et al., “Attention Is All You Need”, \emph{Advances in Neural Information Processing Systems 30 (NeurIPS 2017)}, 2017, pp. 5998—6008.}十分相似：人类作曲家也只有经历大量的音乐积累与训练，才能
实现所谓的“动机”——即便是音乐神童莫扎特，也离不开大量的音乐积累；若让一个连琴键都没
碰过的人立刻完成一部交响乐，那几乎违背学习规律。从理性的角度看，所谓创作“四大件”
——和声、曲式、复调、配器——皆有大量对应的逻辑、书籍与可承载的知识体系，否则整个音乐
构成就会乱成一锅粥；即便是现代派的作曲，也是一种刻意的安排。事实上，二十世纪的序列主义
（如勋伯格的十二音技法、布列兹的总序列主义）已将音乐的可量化推至极端——连音高、时值、
力度、音色都实现了全参数化，而世人并不否认其为音乐。

\subsection{结构的可操作性与“好坏”的归属}

从数学的角度说，这仍然可以解释为“在一定时间内所占据的声音比例”。数字的集合不会因为
你把 0 换成 100 就判定数字消失了；同样，音乐里也不会因为你在某处把一个三和弦替换成七和
弦就彻底失效。更严谨地说：结构操作不会摧毁音乐作为声音组织的有效性，而只改变其审美
效应——审美效应正属于不可量化的运动层。

那么此时，有人一定会问：如何判定好不好呢？这种问题仍然属于“音乐的审美”与“人为的相对
动态”的管辖范围，与音乐本身的结构毫无关联。譬如，不能说贝多芬的《命运》就一定写得比
《粉刷匠》好：从结构上看，《命运》的音乐结构要复杂得多，但结构复杂不等于“好”，而是
取决于听者是谁——对儿童而言，显而易见《粉刷匠》更具启蒙意义，也更容易为儿童接受。
问题的关键仍然回到“音乐的审美”与“人为的相对动态”：因为量化的音乐结构不以人的意志为
转移，所以结构本身无所谓好坏。

\section{量化意识缺位的制度性根源}

\subsection{口传心授与“对错”评价的错位}

既然音乐可量化、且量化意识本应是音乐的内在基因，那么，为什么这个问题长期难以被发现？
从教学方式与社会存在性来看，音乐长期以一种“口传心授”的方式运行，长此以往，“量化”的
念头便在不自觉间消失了。举一个例子：你在这里弹奏，老师说“对的”；你在另一个地方弹奏，
老师说“错了”。但从前面的论证来看，音乐结构本身没有对错，审美层只是千人千面——那么，
“对错”从何而来？可见，教育系统把审美偏好包装成了“对错”来评判与打分，这是一种制度性
的错位：该量化的结构未被量化，不该量化的审美却被量化成了分数。

\subsection{机器裁判的思想实验}

那么，审美层真的可以量化吗？如果真的以所谓“大师”演奏家为参考，最好的方式似乎是让他们
对着 MIDI 键盘录一遍，此后所有考评都拿着学生的演奏与机器比照一番，计算差值。但这显然
会让艺术的表达丢掉自我，且极不可取。值得注意的是，这个思想实验从反面印证了本文的分层：
机器评价“审美”是荒谬的，但机器评价“结构”却完美公正。结论不是“机器不可用”，而是
“机器只能用在该用的层”。

\subsection{学科建制与跨学科思辨的缺失}

从学业与社会性规划来看：音乐长期作为文科的分支，在某种程度上丢失了理科的基因，这致使
跨学科思辨能力减弱。此处须作一说明：并非主张音乐“变成”理科，而是说音乐同时具备可计算
的结构层，这一层需要理科思维的参与——是学科交叉，而非学科替代。其次，当学生把大量时间
放在审美偏好的磨合上时，时间是公平的单位，其在结构层的理解也相应减少，自然也无从萌生
接触额外知识并与音乐结合的可能。

\section{结论与展望：走向分层的音乐量化}

\subsection{四类错位的系统化}

综观全文，“错位”一词可以贯穿始终：本文所揭示的，实际上是四种彼此纠缠的错位，如表
\ref{tab:cuowei} 所示。

\begin{table}[htbp]
  \centering
  \caption{音乐“非量化”工作中的四类错位}
  \label{tab:cuowei}
  \begin{tabular}{p{2.2cm}p{8.4cm}p{3.2cm}}
    \toprule
    \thd{错位类型} & \thd{内涵} & \thd{所在层面} \\
    \midrule
    定义错位 & 将“音乐情感不可量化”混同于“音乐不可量化” & 命题层面 \\
    \addlinespace
    美学错位 & 将审美价值绑定于“人”这一载体，以产生者定义音乐 & 本体层面 \\
    \addlinespace
    技术错位 & 将结构层与审美层混为一谈，或以粗量化杀死音乐，或因畏难放弃量化 & 方法层面 \\
    \addlinespace
    制度错位 & 以“口传心授”传承、以“对错”评价审美、以学科建制隔绝理科思维 & 教育与社会层面 \\
    \bottomrule
  \end{tabular}
\end{table}

\subsection{分层量化方案展望}

基于上述论证，音乐量化不应是“切碎音乐”，而应是“为音乐建立坐标系”。可行的方案是分层
量化：在结构层，以绝对的时间坐标为一切声音事件定位——坐标本身不损失信息，因为它只是
参照；在审美层，让不可量化的呼吸、偏移与微差以“容差参数”的形式在坐标之间显式存活——
偏差不再是错误，而是音乐得以“活着”的证据。

这一方案对 AI 音乐编曲的意义在于：当前音乐 AI 的弱势，不在于缺数据或缺算力，而在于在
错误的坐标系上堆算力。一旦结构层获得精确的绝对坐标、审美层获得显式的容差表达，AI 便
不再需要“模仿”音乐，而可以“书写”音乐。诚然，这仍是未竟的事业，本文不过为它提供一块
本体论的地基。

结语处，笔者愿以一言自况：此事如登高望远，不敢自称登顶，唯愿以寸步之进，换后来者拾级
之便。至于其间锋芒，古人已先我而言——“不恨古人吾不见，恨古人不见吾狂耳”。

\end{document}
